\begin{table}[H]
\centering
\caption{Standardized Effect Sizes for Main Outcomes}
\label{tab:sde}
\begin{threeparttable}
\begin{adjustbox}{max width=\textwidth}
\footnotesize
\begin{tabular}{llcccccc}
\toprule
Outcome & Specification & $\hat{\beta}$ & SE & SD($Y$) & SDE & SE(SDE) & Classification \\
\midrule
\multicolumn{8}{l}{\textit{Panel A: Pooled}} \\
$\Delta$ Occ.\ Score 1930--40 & Preferred IV & 3.167 & 0.911 & 6.239 & 0.508 & 0.146 & Large positive \\
$\Delta$ Occ.\ Score 1920--30 & Boom IV & 3.976 & 1.634 & 6.613 & 0.601 & 0.247 & Large positive \\
\midrule
\multicolumn{8}{l}{\textit{Panel B: Heterogeneous}} \\
$\Delta$ Occ.\ Score 1930--40 & Farm origin & 2.579 & 1.018 & 5.412 & 0.477 & 0.188 & Large positive \\
$\Delta$ Occ.\ Score 1930--40 & Non-farm origin & 3.999 & 1.489 & 7.501 & 0.533 & 0.199 & Large positive \\
\bottomrule
\end{tabular}
\end{adjustbox}
\begin{tablenotes}[flushleft]
\small
\item \textit{Notes:} \textbf{Country:} United States. \textbf{Research question:} Whether the occupational upgrading achieved by Black Great Migration participants during the 1920s boom survived the Great Depression, the first individual-level causal test of the ``last hired, first fired'' hypothesis. \textbf{Policy mechanism:} The Great Migration relocated millions of African Americans from the rural South to Northern industrial cities during 1910--1970; Northern labor markets offered structurally different occupational ladders than Southern agriculture, but the Depression tested whether these gains were durable or whether racial discrimination reversed them. \textbf{Outcome definition:} Change in occupational income score (occscore) between 1930 and 1940 Censuses, where occscore maps each Census occupation code to the median income of workers in that occupation in 1950. \textbf{Treatment:} Binary indicator for South-to-North migration (residing in a Southern state in 1920 and a Northern/Western state by 1930). \textbf{Data:} IPUMS Multigenerational Longitudinal Panel linking individuals across the 1920, 1930, and 1940 Censuses; unit of observation is an individual male aged 15--55 in 1920. \textbf{Method:} Two-stage least squares instrumenting migration with a shift-share measure combining inverse geographic distance from origin counties to 12 Northern cities with 1910 Black population stocks; standard errors clustered at the origin county level. \textbf{Sample:} Black males born in 13 Southern states, aged 15--55 in 1920, with valid occupational scores in all three Census years; excludes those with zero occupational scores. SDE $= \hat{\beta} / \text{SD}(Y)$ where SD($Y$) is the unconditional standard deviation. Classification refers to magnitude, not statistical significance: Large ($|$SDE$| > 0.15$), Moderate ($0.05$--$0.15$), Small ($0.005$--$0.05$), Null ($< 0.005$).
\end{tablenotes}
\end{threeparttable}
\end{table}

